% test-004.tex -- tagged-PDF test de las utilidades para geometria
%
% Compilar:  lualatex-dev test-004.tex
% Validar:   show-pdf-tags --xml test-004.pdf | rnv-wrapp
%            verapdf --flavour ua2 --format text test-004.pdf
%
% Test de prueba para las utilidades para geometría: las formas de cada
% comando y sus llaves. Los "doc:" son lecturas reales de NVDA 2026.2 con MathCAT.
\DocumentMetadata
  {
    lang=es-CL,
    pdfversion=2.0,
    pdfstandard=ua-2,
    tagging=on,
    tagging-setup={math/setup={mathml-SE}},
  }
\documentclass{article}
\usepackage[spanish]{babel}
\usepackage{unicode-math}
\usepackage{spintent}
\usepackage{hyperref}
\hypersetup
  {
    pdfdisplaydoctitle = true,
    pdftitle           = {test-004: utilidades para geometría},
  }

\begin{document}

\section*{A. Los comandos \texttt{\textbackslash spcoord} y \texttt{\textbackslash spPoint}}

% doc: «menos 3 coma 4»
1. Coordenadas: $\spcoord{-3, 4}$.

% doc: «x subíndice 1 coma y subíndice 1»
2. Coordenadas con subíndices: $\spcoord{x_{1},y_{1}}$.

% doc: «(silencio)»
3. Coordenadas silenciadas: $\spcoord[silent=true]{-3, 4}$.

% doc: «punto A»
4. Punto: $\spPoint{A}$.

% doc: «punto A sub uno»
5. Punto con subíndice: $\spPoint{A_{1}}$.

% doc: «punto A menos 3 coma 4»
6. Punto con coordenadas: $\spPoint{A}(-3,4)$.

% doc: «punto A sub uno menos 3 coma 4»
7. Punto con subíndice y coordenadas: $\spPoint{A_{1}}(-3,4)$.

% doc: «punto A 1 coma 5 punto y coma 2»
8. Punto con coordenadas decimales: $\spPoint{A}(1,5;2)$.

% doc: «punto A sub uno»
9. Punto con lectura \texttt{school}: $\spPoint[read-cmd=school]{A_{1}}$.

% doc: «punto a subíndice 1»
10. Punto con lectura \texttt{mathcat}: $\spPoint[read-cmd=mathcat]{A_{1}}$.

% doc: «sea punto A»
11. Punto con prefijo: $\spPoint[prefix={sea}]{A}$.

% doc: «A»
12. Punto sin el concepto: $\spPoint[concept=false]{A}$.

% doc: «(silencio)»
13. Punto silenciado: $\spPoint[silent=true]{A}$.

\section*{B. El comando \texttt{\textbackslash spvec}}

% doc: «vector a»
14. Vector con una letra: $\spvec{a}$.

% doc: «vector A B»
15. Vector con dos puntos: $\spvec{AB}$.

% doc: «vector A B»
16. Vector con dos puntos separados: $\spvec{A,B}$.

% doc: «vector a es igual a menos 3 coma 4»
17. Vector con componentes: $\spvec{a}(-3,4)$.

% doc: «vector A B es igual a menos 3 coma 4»
18. Vector entre puntos con componentes: $\spvec{AB}(-3,4)$.

% doc: «vector A B»
19. Vector con lectura \texttt{school}: $\spvec[read-cmd=school]{AB}$.

% doc: «sea vector a»
20. Vector con prefijo: $\spvec[prefix={sea}]{a}$.

% doc: «vector A B»
21. Vector con ajuste automático: $\spvec[autofit=true]{AB}$.

% doc: «a»
22. Vector sin el concepto: $\spvec[concept=false]{a}$.

% doc: «(silencio)»
23. Vector silenciado: $\spvec[silent=true]{a}$.

\section*{C. Los comandos \texttt{\textbackslash sprecta}, \texttt{\textbackslash sprayo}, \texttt{\textbackslash spside} y \texttt{\textbackslash spmside}}

% doc: «recta A B»
24. Recta por dos puntos: $\sprecta{AB}$.

% doc: «recta A B»
25. Recta por dos puntos separados: $\sprecta{A,B}$.

% doc: «recta m»
26. Recta con una letra: $\sprecta{m}$.

% doc: «recta A B»
27. Recta con lectura \texttt{school}: $\sprecta[read-cmd=school]{AB}$.

% doc: «sea recta A B»
28. Recta con prefijo: $\sprecta[prefix={sea}]{AB}$.

% doc: «A B»
29. Recta sin el concepto: $\sprecta[concept=false]{AB}$.

% doc: «(silencio)»
30. Recta silenciada: $\sprecta[silent=true]{AB}$.

% doc: «rayo A B»
31. Rayo: $\sprayo{AB}$.

% doc: «semirrecta A B»
32. Rayo con lectura \texttt{semirrecta}: $\sprayo[read-sty=semirrecta]{AB}$.

% doc: «rayo A B»
33. Rayo con lectura \texttt{school}: $\sprayo[read-cmd=school]{AB}$.

% doc: «sea rayo A B»
34. Rayo con prefijo: $\sprayo[prefix={sea}]{AB}$.

% doc: «A B»
35. Rayo sin el concepto: $\sprayo[concept=false]{AB}$.

% doc: «(silencio)»
36. Rayo silenciado: $\sprayo[silent=true]{AB}$.

% doc: «lado A B»
37. Lado: $\spside{AB}$.

% doc: «segmento A B»
38. Lado con lectura \texttt{segmento}: $\spside[read-sty=segmento]{AB}$.

% doc: «trazo A B»
39. Lado con lectura \texttt{trazo}: $\spside[read-sty=trazo]{AB}$.

% doc: «lado A B»
40. Lado con lectura \texttt{school}: $\spside[read-cmd=school]{AB}$.

% doc: «sea lado A B»
41. Lado con prefijo: $\spside[prefix={sea}]{AB}$.

% doc: «A B»
42. Lado sin el concepto: $\spside[concept=false]{AB}$.

% doc: «lado base A B»
43. Lado con verbalización propia: $\spside[read-arg={lado base}]{AB}$.

% doc: «(silencio)»
44. Lado silenciado: $\spside[silent=true]{AB}$.

% doc: «medida del lado A B»
45. Medida del lado: $\spmside{AB}$.

% doc: «medida del lado A B»
46. Medida del lado con la letra m: $\spmside[print-sym=true]{AB}$.

% doc: «medida del segmento A B»
47. Medida del lado con lectura \texttt{segmento}: $\spmside[read-sty=segmento]{AB}$.

% doc: «medida del lado A B»
48. Medida del lado con lectura \texttt{school}: $\spmside[read-cmd=school]{AB}$.

% doc: «sea medida del lado A B»
49. Medida del lado con prefijo: $\spmside[prefix={sea}]{AB}$.

% doc: «A B»
50. Medida del lado sin el concepto: $\spmside[concept=false]{AB}$.

% doc: «medida de la base A B»
51. Medida del lado con verbalización propia: $\spmside[concept=false,read-arg={medida de la base}]{AB}$.

% doc: «(silencio)»
52. Medida del lado silenciada: $\spmside[silent=true]{AB}$.

\section*{D. Los comandos \texttt{\textbackslash spangle}, \texttt{\textbackslash spmangle} y \texttt{\textbackslash spang}}

% doc: «ángulo en A»
53. Ángulo en un punto: $\spangle{A}$.

% doc: «ángulo alfa»
54. Ángulo con letra griega: $\spangle{\alpha}$.

% doc: «ángulo 1»
55. Ángulo con número: $\spangle{1}$.

% doc: «ángulo A O B»
56. Ángulo con tres puntos: $\spangle{AOB}$.

% doc: «ángulo A O B»
57. Ángulo con tres puntos separados: $\spangle{A,O,B}$.

% doc: «ángulo recto A O B»
58. Ángulo recto: $\spangle[recto=true]{AOB}$.

% doc: «ángulo A O B»
59. Ángulo con lectura \texttt{school}: $\spangle[read-cmd=school]{AOB}$.

% doc: «sea ángulo A O B»
60. Ángulo con prefijo: $\spangle[prefix={sea}]{AOB}$.

% doc: «A O B»
61. Ángulo sin el concepto: $\spangle[concept=false]{AOB}$.

% doc: «(silencio)»
62. Ángulo silenciado: $\spangle[silent=true]{AOB}$.

% doc: «medida del ángulo en A»
63. Medida del ángulo en un punto: $\spmangle{A}$.

% doc: «medida del ángulo alfa»
64. Medida del ángulo con letra griega: $\spmangle{\alpha}$.

% doc: «medida del ángulo 1»
65. Medida del ángulo con número: $\spmangle{1}$.

% doc: «medida del ángulo A O B»
66. Medida del ángulo con tres puntos: $\spmangle{AOB}$.

% doc: «medida del ángulo recto A O B»
67. Medida del ángulo recto: $\spmangle[recto=true]{AOB}$.

% doc: «medida del ángulo A O B»
68. Medida del ángulo con la letra m: $\spmangle[print-sym=true]{AOB}$.

% doc: «medida del ángulo A O B»
69. Medida del ángulo con lectura \texttt{school}: $\spmangle[read-cmd=school]{AOB}$.

% doc: «sea medida del ángulo A O B»
70. Medida del ángulo con prefijo: $\spmangle[prefix={sea}]{AOB}$.

% doc: «A O B»
71. Medida del ángulo sin el concepto: $\spmangle[concept=false]{AOB}$.

% doc: «(silencio)»
72. Medida del ángulo silenciada: $\spmangle[silent=true]{AOB}$.

% doc: «45 grados»
73. Grados: $\spang{45}$.

% doc: «45 grados 30 minutos»
74. Grados y minutos: $\spang{45:30}$.

% doc: «45 grados 30 minutos 15 segundos»
75. Grados, minutos y segundos: $\spang{45:30:15}$.

\section*{E. Los comandos \texttt{\textbackslash sparc} y \texttt{\textbackslash spmarc}}

% doc: «arco A B»
76. Arco: $\sparc{AB}$.

% doc: «arco A B C»
77. Arco con tres puntos: $\sparc{ABC}$.

% doc: «arco A B»
78. Arco con lectura \texttt{school}: $\sparc[read-cmd=school]{AB}$.

% doc: «sea arco A B»
79. Arco con prefijo: $\sparc[prefix={sea}]{AB}$.

% doc: «A B»
80. Arco sin el concepto: $\sparc[concept=false]{AB}$.

% doc: «(silencio)»
81. Arco silenciado: $\sparc[silent=true]{AB}$.

% doc: «medida del arco A B»
82. Medida del arco: $\spmarc{AB}$.

% doc: «medida del arco A B C»
83. Medida del arco con tres puntos: $\spmarc{ABC}$.

% doc: «medida del arco A B»
84. Medida del arco con lectura \texttt{school}: $\spmarc[read-cmd=school]{AB}$.

% doc: «sea medida del arco A B»
85. Medida del arco con prefijo: $\spmarc[prefix={sea}]{AB}$.

% doc: «A B»
86. Medida del arco sin el concepto: $\spmarc[concept=false]{AB}$.

% doc: «(silencio)»
87. Medida del arco silenciada: $\spmarc[silent=true]{AB}$.

\section*{F. El comando \texttt{\textbackslash spPoly}}

% doc: «triángulo A B C»
88. Triángulo: $\spPoly{ABC}$.

% doc: «cuadrado A B C D»
89. Cuadrilátero: $\spPoly{ABCD}$.

% doc: «pentágono A B C D E»
90. Pentágono: $\spPoly{ABCDE}$.

% doc: «triángulo A B C»
91. Triángulo sin símbolo: $\spPoly[print-sym=false]{ABC}$.

% doc: «cuadrado A B C D»
92. Cuadrilátero sin símbolo: $\spPoly[print-sym=false]{ABCD}$.

% doc: «triángulo A B C»
93. Polígono con lectura \texttt{school}: $\spPoly[read-cmd=school]{ABC}$.

% doc: «sea triángulo A B C»
94. Polígono con prefijo: $\spPoly[prefix={sea}]{ABC}$.

% doc: «A B C»
95. Polígono sin el concepto: $\spPoly[concept=false]{ABC}$.

% doc: «(silencio)»
96. Polígono silenciado: $\spPoly[silent=true]{ABC}$.

\section*{G. Los comandos \texttt{\textbackslash spAfig} y \texttt{\textbackslash spPfig}}

% doc: «área»
97. Área sin argumento: $\spAfig$.

% doc: «área uno»
98. Área de la figura número uno: $\spAfig(1)$.

% doc: «área n»
99. Área de la figura número \emph{n}: $\spAfig(n)$.

% doc: «área del triángulo A B C»
100. Área de un triángulo: $\spAfig(ABC)$.

% doc: «área del cuadrado A B C D»
101. Área de un cuadrilátero: $\spAfig(ABCD)$.

% doc: «área del pentágono A B C D E»
102. Área de un pentágono: $\spAfig(ABCDE)$.

% doc: «área de la figura»
103. Área de una figura: $\spAfig(F)$.

% doc: «área de la figura uno»
104. Área de una figura con subíndice: $\spAfig(F_{1})$.

% doc: «área del polígono»
105. Área de un polígono: $\spAfig(P)$.

% doc: «área basal»
106. Área basal: $\spAfig(B)$.

% doc: «área lateral»
107. Área lateral: $\spAfig(L)$.

% doc: «área total»
108. Área total: $\spAfig(T)$.

% doc: «área del trapecio A B C D»
109. Área con nombre propio: $\spAfig[read-arg={trapecio}](ABCD)$.

% doc: «área sub uno»
110. Área con subíndice leído: $\spAfig[read-sub=true](1)$.

% doc: «área del triángulo A B C»
111. Área sin símbolo: $\spAfig[print-sym=none](ABC)$.

% doc: «área del triángulo uno»
112. Área con símbolo de triángulo: $\spAfig[print-sym=triangle](1)$.

% doc: «área del cuadrado uno»
113. Área con símbolo de cuadrado: $\spAfig[print-sym=square](1)$.

% doc: «área del triángulo A B C»
114. Área con lectura \texttt{school}: $\spAfig[read-cmd=school](ABC)$.

% doc: «calcule área del triángulo A B C»
115. Área con prefijo: $\spAfig[prefix={calcule}](ABC)$.

% doc: «A B C»
116. Área sin el concepto: $\spAfig[concept=false](ABC)$.

% doc: «(silencio)»
117. Área silenciada: $\spAfig[silent=true](ABC)$.

% doc: «perímetro»
118. Perímetro sin argumento: $\spPfig$.

% doc: «perímetro uno»
119. Perímetro de la figura número uno: $\spPfig(1)$.

% doc: «perímetro n»
120. Perímetro de la figura número \emph{n}: $\spPfig(n)$.

% doc: «perímetro del triángulo A B C»
121. Perímetro de un triángulo: $\spPfig(ABC)$.

% doc: «perímetro del cuadrado A B C D»
122. Perímetro de un cuadrilátero: $\spPfig(ABCD)$.

% doc: «perímetro del pentágono A B C D E»
123. Perímetro de un pentágono: $\spPfig(ABCDE)$.

% doc: «perímetro de la figura»
124. Perímetro de una figura: $\spPfig(F)$.

% doc: «perímetro de la figura uno»
125. Perímetro de una figura con subíndice: $\spPfig(F_{1})$.

% doc: «perímetro del polígono»
126. Perímetro de un polígono: $\spPfig(P)$.

% doc: «perímetro basal»
127. Perímetro basal: $\spPfig(B)$.

% doc: «perímetro lateral»
128. Perímetro lateral: $\spPfig(L)$.

% doc: «perímetro total»
129. Perímetro total: $\spPfig(T)$.

% doc: «perímetro del trapecio A B C D»
130. Perímetro con nombre propio: $\spPfig[read-arg={trapecio}](ABCD)$.

% doc: «perímetro sub uno»
131. Perímetro con subíndice leído: $\spPfig[read-sub=true](1)$.

% doc: «perímetro del triángulo A B C»
132. Perímetro sin símbolo: $\spPfig[print-sym=none](ABC)$.

% doc: «perímetro del triángulo uno»
133. Perímetro con símbolo de triángulo: $\spPfig[print-sym=triangle](1)$.

% doc: «perímetro del cuadrado uno»
134. Perímetro con símbolo de cuadrado: $\spPfig[print-sym=square](1)$.

% doc: «perímetro del triángulo A B C»
135. Perímetro con lectura \texttt{school}: $\spPfig[read-cmd=school](ABC)$.

% doc: «calcule perímetro del triángulo A B C»
136. Perímetro con prefijo: $\spPfig[prefix={calcule}](ABC)$.

% doc: «A B C»
137. Perímetro sin el concepto: $\spPfig[concept=false](ABC)$.

% doc: «(silencio)»
138. Perímetro silenciado: $\spPfig[silent=true](ABC)$.

\section*{H. Los comandos \texttt{\textbackslash sprad} y \texttt{\textbackslash spdiam}}

% doc: «radio»
139. Radio: $\sprad$.

% doc: «diámetro»
140. Diámetro: $\spdiam$.

% doc: «radio»
141. Radio en mayúscula: $\sprad[upper-case=true]$.

% doc: «diámetro»
142. Diámetro en mayúscula: $\spdiam[upper-case=true]$.

% doc: «diámetro»
143. Diámetro con símbolo: $\spdiam[print-sym=true]$.

% doc: «radio 1»
144. Radio con subíndice: $\sprad[sub-index={1}]$.

% doc: «radio sub 1»
145. Radio con subíndice leído: $\sprad[sub-index={1},read-sub=true]$.

% doc: «radio»
146. Radio con lectura \texttt{school}: $\sprad[read-cmd=school]$.

% doc: «r»
147. Radio sin el concepto: $\sprad[concept=false]$.

% doc: «medida del radio»
148. Radio con verbalización propia: $\sprad[concept=false,read-arg={medida del radio}]$.

% doc: «(silencio)»
149. Radio silenciado: $\sprad[silent=true]$.

\section*{I. Los comandos \texttt{\textbackslash spcirc}, \texttt{\textbackslash spAcirc} y \texttt{\textbackslash spPcirc}}

% doc: «circunferencia»
150. Circunferencia: $\spcirc$.

% doc: «circunferencia de centro en O»
151. Circunferencia con centro: $\spcirc(O)$.

% doc: «circunferencia de centro en O y radio 3»
152. Circunferencia con centro y radio: $\spcirc(O,3)$.

% doc: «circunferencia de centro en O y radio medida del lado A B»
153. Circunferencia con radio dado por un segmento: $\spcirc(O, AB)$.

% doc: «circunferencia de centro en O y radio lado A B»
154. Circunferencia con radio dado por un segmento (\texttt{side}): $\spcirc[radio-sty=side](O, AB)$.

% doc: «circunferencia de centro en O y radio r»
155. Circunferencia con radio dado por una letra: $\spcirc(O, r)$.

% doc: «circunferencia de centro en O y radio 3 centímetros»
156. Circunferencia con radio como medida: $\spcirc(O,3 cm)$.

% doc: «circunferencia de centro en O sub uno»
157. Circunferencia con centro con subíndice: $\spcirc(O_{1})$.

% doc: «circunferencia de centro en O»
158. Circunferencia con símbolo usual: $\spcirc[print-sym=usual](O)$.

% doc: «circunferencia de centro en O y diámetro 3»
159. Circunferencia con diámetro: $\spcirc[diameter=true](O,3)$.

% doc: «círculo de centro en O»
160. Circunferencia con lectura \texttt{circulo}: $\spcirc[read-sty=circulo](O)$.

% doc: «circunferencia 1 de centro en O y radio medida del lado A B»
161. Circunferencia con subíndice: $\spcirc[sub-index={1}](O, AB)$.

% doc: «circunferencia de centro en O»
162. Circunferencia con lectura \texttt{school}: $\spcirc[read-cmd=school](O)$.

% doc: «sea circunferencia de centro en O»
163. Circunferencia con prefijo: $\spcirc[prefix={sea}](O)$.

% doc: «O»
164. Circunferencia sin el concepto: $\spcirc[concept=false](O)$.

% doc: «circunferencia unidad de centro en O»
165. Circunferencia con verbalización propia: $\spcirc[read-arg={circunferencia unidad}](O)$.

% doc: «(silencio)»
166. Circunferencia silenciada: $\spcirc[silent=true](O)$.

% doc: «área de la circunferencia»
167. Área del círculo: $\spAcirc$.

% doc: «área de la circunferencia con centro en O»
168. Área del círculo con centro: $\spAcirc(O)$.

% doc: «área de la circunferencia uno»
169. Área del círculo número uno: $\spAcirc(1)$.

% doc: «área de la circunferencia con centro en O»
170. Área del círculo con símbolo usual: $\spAcirc[print-sym=usual](O)$.

% doc: «área de la circunferencia 1 con centro en O»
171. Área del círculo con subíndice: $\spAcirc[sub-index={1}](O)$.

% doc: «área del círculo con centro en O»
172. Área del círculo con lectura \texttt{circulo}: $\spAcirc[read-sty=circulo](O)$.

% doc: «perímetro de la circunferencia»
173. Perímetro de la circunferencia: $\spPcirc$.

% doc: «perímetro de la circunferencia con centro en O»
174. Perímetro de la circunferencia con centro: $\spPcirc(O)$.

% doc: «perímetro de la circunferencia uno»
175. Perímetro de la circunferencia número uno: $\spPcirc(1)$.

% doc: «perímetro de la circunferencia con centro en O»
176. Perímetro de la circunferencia con símbolo usual: $\spPcirc[print-sym=usual](O)$.

% doc: «perímetro de la circunferencia 1 con centro en O»
177. Perímetro de la circunferencia con subíndice: $\spPcirc[sub-index={1}](O)$.

% doc: «perímetro del círculo con centro en O»
178. Perímetro de la circunferencia con lectura \texttt{circulo}: $\spPcirc[read-sty=circulo](O)$.

\section*{J. El comando \texttt{\textbackslash spVol}}

% doc: «volumen»
179. Volumen: $\spVol$.

% doc: «volumen 1»
180. Volumen con subíndice: $\spVol[sub-index={1}]$.

% doc: «volumen sub 1»
181. Volumen con subíndice leído: $\spVol[sub-index={1},read-sub=true]$.

% doc: «volumen»
182. Volumen con lectura \texttt{school}: $\spVol[read-cmd=school]$.

% doc: «V»
183. Volumen sin el concepto: $\spVol[concept=false]$.

% doc: «volumen del cilindro»
184. Volumen con verbalización propia: $\spVol[concept=false,read-arg={volumen del cilindro}]$.

% doc: «(silencio)»
185. Volumen silenciado: $\spVol[silent=true]$.

\end{document}
